\documentclass[11pt,a4paper]{article}
\usepackage[T1]{fontenc}
\usepackage[{{babel}}]{babel}
\usepackage{lmodern}
\usepackage[margin=2.5cm]{geometry}
\usepackage{amsmath}
\usepackage{graphicx}
\usepackage{booktabs}
\usepackage{siunitx}
\sisetup{locale = {{fr:FR|en:US}}}
\usepackage{pgfplots}
\pgfplotsset{compat=1.18}
\usepackage{float}
\usepackage[hidelinks]{hyperref}

\begin{document}

% {{fr:Pas d'ancre de page pour la page de titre (évite un doublon « page.1 »)|en:No page anchor for the title page (avoids a duplicate "page.1")}}
\hypersetup{pageanchor=false}
\begin{titlepage}
  \centering
  {\Large {{institution}}\par}
  \vspace{3cm}
  {\huge\bfseries {{title}}\par}
  \vspace{1cm}
  {\Large {{fr:Compte rendu de travaux pratiques|en:Lab report}}\par}
  \vfill
  {\large {{author}}\par}
  {{date}}
\end{titlepage}
\hypersetup{pageanchor=true}

\section{{{fr:Objectifs|en:Objectives}}}
{{fr:Vérifier la loi d'Ohm $U = R I$ et mesurer la résistance $R$.|en:Check Ohm's law $U = R I$ and measure the resistance $R$.}}

\section{{{fr:Matériel et protocole|en:Equipment and method}}}
\begin{itemize}
  \item {{fr:Générateur de tension continue (\qtyrange{0}{12}{\volt})|en:DC voltage supply (\qtyrange{0}{12}{\volt})}}
  \item {{fr:Résistance nominale de \qty{100}{\ohm}|en:Resistor of nominal value \qty{100}{\ohm}}}
  \item {{fr:Deux multimètres|en:Two multimeters}}
\end{itemize}

\section{{{fr:Résultats|en:Results}}}

\begin{table}[H]
  \centering
  \caption{{{fr:Mesures|en:Measurements}}}
  \label{tab:mesures}
  \begin{tabular}{S[table-format=2.1] S[table-format=3.1]}
    \toprule
    {$U$ (\si{\volt})} & {$I$ (\si{\milli\ampere})} \\
    \midrule
    2.0  & 20.1 \\
    4.0  & 39.8 \\
    6.0  & 60.3 \\
    8.0  & 79.9 \\
    10.0 & 100.2 \\
    \bottomrule
  \end{tabular}
\end{table}

\begin{figure}[H]
  \centering
  \begin{tikzpicture}
    \begin{axis}[
      width=0.75\linewidth, height=6cm,
      xlabel={$I$ (\si{\milli\ampere})}, ylabel={$U$ (\si{\volt})},
      grid=major, legend pos=north west]
      \addplot[only marks, mark=*, blue] coordinates {(20.1,2) (39.8,4) (60.3,6) (79.9,8) (100.2,10)};
      \addlegendentry{{{fr:Mesures|en:Measurements}}}
      \addplot[red, domain=0:110] {0.0998*x};
      \addlegendentry{{{fr:Régression|en:Fit}} $U = RI$}
    \end{axis}
  \end{tikzpicture}
  \caption{{{fr:Caractéristique $U(I)$ de la résistance.|en:$U(I)$ characteristic of the resistor.}}}
  \label{fig:caracteristique}
\end{figure}

{{fr:La pente de la droite (figure~\ref{fig:caracteristique}) donne $R = \qty{99.8 \pm 0.5}{\ohm}$.|en:The slope of the line (Figure~\ref{fig:caracteristique}) gives $R = \qty{99.8 \pm 0.5}{\ohm}$.}}

\section{Discussion}
{{fr:Comparer avec la valeur nominale et discuter des incertitudes.|en:Compare with the nominal value and discuss uncertainties.}}

\end{document}
